\chapter{2004年北京卷（春理）}
\section{单选题}
\begin{problem}
在函数 \(y=\sin 2x\)，\(y=\sin x\)，\(y=\cos x\)，\(y=\tan \dfrac{x}{2}\) 中，最小正周期为 \(\pi\) 的函数是（\quad）

\choices
    {\(y=\sin 2x\)}
    {\(y=\sin x\)}
    {\(y=\cos x\)}
    {\(y=\tan \dfrac{x}{2}\)}
\end{problem}

\begin{answer}
A
\end{answer}

\begin{solution}
函数 \(y=\sin kx\) 或 \(y=\cos kx\) 的最小正周期为 \(\dfrac{2\pi}{|k|}\). 因此，\(y=\sin 2x\) 的最小正周期为 \(\dfrac{2\pi}{2}=\pi\)，而 \(y=\sin x\)、\(y=\cos x\) 的最小正周期均为 \(2\pi\). 对于 \(y=\tan\dfrac{x}{2}\)，正切函数的最小正周期为 \(\pi\)，所以它的最小正周期为 \(2\pi\)，不是 \(\pi\). 故选 A.
\end{solution}

\begin{problem}
当 \(\dfrac{2}{3}<m<1\) 时，复数 \(z=(3m-2)+(m-1)\i\) 在复平面上对应的点位于（\quad）

\choices
    {第一象限}
    {第二象限}
    {第三象限}
    {第四象限}
\end{problem}

\begin{answer}
D
\end{answer}

\begin{solution}
由 \(\dfrac{2}{3}<m<1\)，得 \(3m-2>0\)，且 \(m-1<0\). 复数 \(z=(3m-2)+(m-1)\i\) 的实部为正、虚部为负，所以它在复平面上对应的点位于第四象限，故选 D.
\end{solution}

\begin{problem}
双曲线 \(\dfrac{x^2}{4}-\dfrac{y^2}{9}=1\) 的渐近线方程是（\quad）

\choices
    {\(y=\pm\dfrac{3}{2}x\)}
    {\(y=\pm\dfrac{2}{3}x\)}
    {\(y=\pm\dfrac{9}{4}x\)}
    {\(y=\pm\dfrac{4}{9}x\)}
\end{problem}

\begin{answer}
A
\end{answer}

\begin{solution}
将双曲线方程中的常数项移去，得渐近线满足 \(\dfrac{x^2}{4}-\dfrac{y^2}{9}=0\)，即 \(y^2=\dfrac{9}{4}x^2\). 因而 \(y=\pm\dfrac{3}{2}x\)，故选 A.
\end{solution}

\begin{problem}
一个圆锥的侧面积是其底面积的 \(2\text{ 倍}\)，则该圆锥的母线与底面所成的角为（\quad）

\choices
    {\(30^{\circ}\)}
    {\(45^{\circ}\)}
    {\(60^{\circ}\)}
    {\(75^{\circ}\)}
\end{problem}

\begin{answer}
C
\end{answer}

\begin{solution}
设圆锥的底面半径为 \(r\)，母线长为 \(l\). 由侧面积是底面积的 \(2\) 倍，得 \(\pi rl=2\pi r^2\)，所以 \(l=2r\). 设母线与底面所成的角为 \(\alpha\)，在圆锥的轴截面中，母线在底面上的射影长为 \(r\)，故 \(\cos\alpha=\dfrac{r}{l}=\dfrac12\). 因为 \(0^{\circ}<\alpha<90^{\circ}\)，所以 \(\alpha=60^{\circ}\)，故选 C.
\end{solution}

\begin{problem}
在极坐标系中，圆心在 \((\sqrt{2},\pi)\) 且过极点的圆的方程为（\quad）

\choices
    {\(\rho=2\sqrt{2}\cos\theta\)}
    {\(\rho=-2\sqrt{2}\cos\theta\)}
    {\(\rho=2\sqrt{2}\sin\theta\)}
    {\(\rho=-2\sqrt{2}\sin\theta\)}
\end{problem}

\begin{answer}
B
\end{answer}

\begin{solution}
极点为圆上的一点，圆心的直角坐标为 \((-\sqrt{2},0)\)，所以圆的半径为 \(\sqrt{2}\)，其直角坐标方程为 \((x+\sqrt{2})^2+y^2=2\). 代入 \(x=\rho\cos\theta\)、\(y=\rho\sin\theta\)，得 \(\rho^2+2\sqrt{2}\rho\cos\theta=0\). 除去表示极点的因子 \(\rho\) 后，圆的极坐标方程为 \(\rho=-2\sqrt{2}\cos\theta\)，故选 B.
\end{solution}

\begin{problem}
已知 \(\sin(\theta+\pi)<0\)，\(\cos(\theta-\pi)>0\)，则下列不等关系中必定成立的是（\quad）

\choices
    {\(\tan\dfrac{\theta}{2}<\cot\dfrac{\theta}{2}\)}
    {\(\tan\dfrac{\theta}{2}>\cot\dfrac{\theta}{2}\)}
    {\(\sin\dfrac{\theta}{2}<\cos\dfrac{\theta}{2}\)}
    {\(\sin\dfrac{\theta}{2}>\cos\dfrac{\theta}{2}\)}
\end{problem}

\begin{answer}
B
\end{answer}

\begin{solution}
由 \(\sin(\theta+\pi)=-\sin\theta<0\)，得 \(\sin\theta>0\)；由 \(\cos(\theta-\pi)=-\cos\theta>0\)，得 \(\cos\theta<0\). 因而存在 \(k\in\Z\)，使得 \(\dfrac{\pi}{2}+2k\pi<\theta<\pi+2k\pi\). 于是 \(\dfrac{\theta}{2}-k\pi\in\left(\dfrac{\pi}{4},\dfrac{\pi}{2}\right)\)，从而 \(\tan\dfrac{\theta}{2}=\tan\left(\dfrac{\theta}{2}-k\pi\right)>1>\cot\left(\dfrac{\theta}{2}-k\pi\right)=\cot\dfrac{\theta}{2}\)，故选 B.
\end{solution}

\begin{problem}
已知三个不等式：\(ab>0\)，\(bc-ad>0\)，\(\dfrac{c}{a}-\dfrac{d}{b}>0\)（其中 \(a\)，\(b\)，\(c\)，\(d\) 均为实数），用其中两个不等式作为条件，余下的一个不等式作为结论组成一个命题，可组成的正确命题的个数是（\quad）

\choices
    {0}
    {1}
    {2}
    {3}
\end{problem}

\begin{answer}
D
\end{answer}

\begin{solution}
因为不等式 \(ab>0\) 成立，所以 \(a\ne0\)，\(b\ne0\)，并且 \(\dfrac{c}{a}-\dfrac{d}{b}=\dfrac{bc-ad}{ab}\). 记三个不等式分别为 \(A\)、\(B\)、\(C\)，则 \(A\) 表示 \(ab>0\)，\(B\) 表示 \(bc-ad>0\)，\(C\) 表示 \(\dfrac{bc-ad}{ab}>0\).

若以 \(A\)、\(B\) 为条件，则分子、分母均为正，推出 \(C\)；若以 \(A\)、\(C\) 为条件，则分母为正且商为正，推出分子为正，即 \(B\)；若以 \(B\)、\(C\) 为条件，则分子为正且商为正，推出分母为正，即 \(A\). 所以三种选取方式均可组成正确命题，正确命题的个数为 \(3\)，故选 D.
\end{solution}

\begin{problem}
两个完全相同的长方体的长，宽，高分别为 \(5\mathrm{~cm}\)，\(4\mathrm{~cm}\)，\(3\mathrm{~cm}\)，把它们重叠在一起组成一个新长方体，在这些新长方体中，最长的对角线的长度是（\quad）

\choices
    {\(\sqrt{77}\mathrm{~cm}\)}
    {\(7\sqrt{2}\mathrm{~cm}\)}
    {\(5\sqrt{5}\mathrm{~cm}\)}
    {\(10\sqrt{2}\mathrm{~cm}\)}
\end{problem}

\begin{answer}
C
\end{answer}

\begin{solution}
两个长方体可以分别沿面积为 \(4\times3\)、\(5\times3\) 或 \(5\times4\) 的全等面拼接. 对应的新长方体的三条棱长分别为 \((10,4,3)\)、\((5,8,3)\) 和 \((5,4,6)\)，其对角线长度的平方分别为 \(10^2+4^2+3^2=125\)、\(5^2+8^2+3^2=98\)、\(5^2+4^2+6^2=77\). 因为 \(125>98>77\)，所以最长的对角线长度为 \(\sqrt{125}=5\sqrt{5}\,\mathrm{cm}\)，故选 C.
\end{solution}

\begin{problem}
在 \(100\text{ 件}\) 产品中有 \(6\text{ 件}\) 次品，现从中任取 \(3\text{ 件}\) 产品，至少有 \(1\text{ 件}\) 次品的不同取法的种数是（\quad）

\choices
    {\(\mathrm C_6^1\mathrm C_{94}^2\)}
    {\(\mathrm C_6^1\mathrm C_{99}^2\)}
    {\(\mathrm C_{100}^3-\mathrm C_{94}^3\)}
    {\(\mathrm A_{100}^3-\mathrm A_{94}^3\)}
\end{problem}

\begin{answer}
C
\end{answer}

\begin{solution}
从 \(100\) 件产品中任取 \(3\) 件的取法总数为 \(\mathrm C_{100}^{3}\). “至少有 \(1\) 件次品”的反面是“\(3\) 件全为正品”，而正品有 \(100-6=94\) 件，全为正品的取法数为 \(\mathrm C_{94}^{3}\). 因此至少有 \(1\) 件次品的取法数为 \(\mathrm C_{100}^{3}-\mathrm C_{94}^{3}\)，故选 C.
\end{solution}

\begin{problem}
期中考试以后，班长算出了全班 \(40\text{ 人}\) 数学成绩的平均分为 \(M\)，如果把 \(M\) 当成一个同学的分数，与原来的 \(40\text{ 个}\) 分数一起，算出这 \(41\text{ 个}\) 分数的平均值为 \(N\)，那么 \(M:N\) 为（\quad）

\choices
    {\(\dfrac{40}{41}\)}
    {\(1\)}
    {\(\dfrac{41}{40}\)}
    {\(2\)}
\end{problem}

\begin{answer}
B
\end{answer}

\begin{solution}
设全班 \(40\) 人的数学成绩总和为 \(S\). 由平均分为 \(M\)，得 \(S=40M\). 把 \(M\) 作为一个新分数加入后，\(41\) 个分数的总和为 \(40M+M=41M\)，所以新的平均值 \(N=\dfrac{41M}{41}=M\). 因而 \(M:N=1:1=1\)，故选 B.
\end{solution}

\section{填空题}
\begin{problem}
若 \(f^{-1}(x)\) 为函数 \(f(x)=\lg(x+1)\) 的反函数，则 \(f^{-1}(x)\) 的值域是\fillinblank{}.
\end{problem}

\begin{answer}
\(\left(-1,+\infty\right)\)
\end{answer}

\begin{solution}
函数 \(f(x)=\lg(x+1)\) 的定义域为 \(x>-1\)，值域为 \(\R\). 反函数 \(f^{-1}(x)\) 的值域等于原函数 \(f(x)\) 的定义域，所以 \(f^{-1}(x)\) 的值域为 \(\left(-1,+\infty\right)\).
\end{solution}

\begin{problem}
\(\dfrac{\sin(\alpha+30^{\circ})-\sin(\alpha-30^{\circ})}{\cos\alpha}\) 的值为\fillinblank{}.
\end{problem}

\begin{answer}
\(1\)
\end{answer}

\begin{solution}
在原式有意义的条件下，\(\cos\alpha\ne0\). 由正弦差公式，得 \(\sin(\alpha+30^{\circ})-\sin(\alpha-30^{\circ})=2\cos\alpha\sin30^{\circ}=\cos\alpha\)，所以原式等于 \(1\).
\end{solution}

\begin{problem}
据某校环保小组调查，某区垃圾量的年增长率为 \(b\)，\(2003\text{ 年}\) 产生的垃圾量为 \(a\text{ 吨}\). 由此预测，该区下一年的垃圾量为\fillinblank{}\(\text{吨}\)，\(2008\text{ 年}\) 的垃圾量为\fillinblank{}\(\text{吨}\).
\end{problem}

\begin{answer}
\(a(1+b)\)；\(a(1+b)^5\)
\end{answer}

\begin{solution}
年增长率为 \(b\)，表示下一年的数量是本年的 \(1+b\) 倍. 因此 \(2003\) 年后的下一年，即 \(2004\) 年的垃圾量为 \(a(1+b)\text{ 吨}\). \(2008\) 年比 \(2003\) 年晚 \(5\) 年，连续增长 \(5\) 次，所以垃圾量为 \(a(1+b)^5\text{ 吨}\).
\end{solution}

\begin{problem}
若直线 \(mx+ny-3=0\) 与圆 \(x^2+y^2=3\) 没有公共点，则 \(m\)，\(n\) 满足的关系式为\fillinblank{}；以 \((m,n)\) 为点 \(P\) 的坐标，过点 \(P\) 的一条直线与椭圆 \(\dfrac{x^2}{7}+\dfrac{y^2}{3}=1\) 的公共点有\fillinblank{}\(\text{个}\).
\end{problem}

\begin{answer}
\(m^2+n^2<3\)；\(2\)
\end{answer}

\begin{solution}
若 \(m=n=0\)，则方程 \(mx+ny-3=0\) 化为 \(-3=0\)，没有点满足，且 \(m^2+n^2=0<3\). 若 \((m,n)\ne(0,0)\)，则直线 \(mx+ny-3=0\) 到原点的距离为 \(\dfrac{3}{\sqrt{m^2+n^2}}\). 该直线与半径为 \(\sqrt{3}\) 的圆没有公共点，当且仅当这段距离大于圆的半径，因此 \(\dfrac{3}{\sqrt{m^2+n^2}}>\sqrt{3}\)，从而 \(m^2+n^2<3\). 综上，所求关系式为 \(m^2+n^2<3\).

由 \(m^2+n^2<3\)，有 \(\dfrac{m^2}{7}+\dfrac{n^2}{3}\le\dfrac{m^2+n^2}{3}<1\)，所以点 \(P(m,n)\) 在椭圆 \(\dfrac{x^2}{7}+\dfrac{y^2}{3}=1\) 的内部. 任取一条过 \(P\) 的直线，设其方向向量为 \((u,v)\ne(0,0)\)，则直线上的点可表示为 \((x,y)=(m,n)+t(u,v)\). 代入椭圆方程，令左边减去 \(1\) 所得的式子为 \(F(t)\)，则 \(F(t)=0\) 是关于 \(t\) 的二次方程，其二次项系数 \(A=\dfrac{u^2}{7}+\dfrac{v^2}{3}>0\)，且 \(C=F(0)=\dfrac{m^2}{7}+\dfrac{n^2}{3}-1<0\). 将其写成 \(At^2+2Ht+C=0\)，其判别式为 \(4(H^2-AC)>0\)，故有两个不相等的实根；两根之积为 \(\dfrac{C}{A}<0\)，所以两根异号. 因此该直线与椭圆有且只有 \(2\) 个公共点.
\end{solution}

\section{解答题}
\begin{problem}
当 \(0<a<1\) 时，解关于 \(x\) 的不等式：\(a^{\sqrt{2x-1}}<a^{x-2}\).
\end{problem}

\begin{solution}
由根式有意义，得 \(2x-1\ge0\)，即 \(x\ge\dfrac12\). 因为 \(0<a<1\)，函数 \(y=a^t\) 在 \(\R\) 上单调递减，所以原不等式等价于 \(\sqrt{2x-1}>x-2\). 令 \(t=\sqrt{2x-1}\)，则 \(t\ge0\)，且 \(x=\dfrac{t^2+1}{2}\)，于是 \(t>\dfrac{t^2+1}{2}-2\)，即 \(t^2-2t-3<0\)，所以 \(\left(t-3\right)\left(t+1\right)<0\). 结合 \(t\ge0\)，得 \(0\le t<3\)，从而 \(0\le2x-1<9\). 因此不等式的解集为 \(\left[\dfrac12,5\right)\).
\end{solution}

\begin{problem}
在 \(\triangle ABC\) 中，\(a\)，\(b\)，\(c\) 分别是 \(\angle A\)，\(\angle B\)，\(\angle C\) 的对边长，已知 \(a\)，\(b\)，\(c\) 成等比数列，且 \(a^2-c^2=ac-bc\)，求 \(\angle A\) 的大小及 \(\dfrac{b\sin B}{c}\) 的值.
\end{problem}

\begin{solution}
因为 \(a\)，\(b\)，\(c\) 是三角形的边长，均为正数；又因为它们成等比数列，所以 \(b^2=ac\). 令 \(q=\dfrac{a}{b}>0\)，则 \(a=qb\)，\(c=\dfrac{b}{q}\). 将其代入已知等式，得 \(b^2\left(q^2-\dfrac1{q^2}\right)=b^2\left(1-\dfrac1q\right)\). 约去 \(b^2\)，再乘以 \(q^2\)，得 \(q^4-1=q^2-q\)，即 \(\left(q-1\right)\left(q^3+q^2+1\right)=0\). 由于 \(q>0\)，有 \(q^3+q^2+1>0\)，故 \(q=1\)，从而 \(a=b=c\). 所以 \(\triangle ABC\) 为等边三角形，\(\angle A=\angle B=60^{\circ}\)，且 \(\dfrac{b\sin B}{c}=\sin60^{\circ}=\dfrac{\sqrt{3}}2\).
\end{solution}

\begin{problem}
如图，四棱锥 \(S-ABCD\) 的底面是边长为 \(1\) 的正方形，\(SD\) 垂直于底面 \(ABCD\)，\(SB=\sqrt{3}\). 求：

\begin{enumerate}
\item 证明：\(BC\perp SC\)；
\item 求面 \(ASD\) 与面 \(BSC\) 所成二面角的大小；
\item 设棱 \(SA\) 的中点为 \(M\)，求异面直线 \(DM\) 与 \(SB\) 所成的角的大小.
\end{enumerate}

\bitmapinlinefigure[6.0cm][]{img/2004/beijing_science_spring/q17_fig1.png}
\end{problem}

\begin{solution}
\begin{enumerate}
\item 建立空间直角坐标系，取 \(D\) 为原点，令底面为 \(z=0\)，并设 \(C(1,0,0)\)，\(A(0,1,0)\)，\(B(1,1,0)\). 设 \(S(0,0,h)\)，其中 \(h>0\). 由 \(SB=\sqrt3\)，得 \(1^2+1^2+h^2=3\)，所以 \(h=1\)，即 \(S(0,0,1)\). 于是 \(\overrightarrow{BC}=(0,-1,0)\)，\(\overrightarrow{SC}=(1,0,-1)\)，从而 \(\overrightarrow{BC}\cdot\overrightarrow{SC}=0\)，故 \(BC\perp SC\).
\item 平面 \(ASD\) 的方程为 \(x=0\)，可取其法向量 \(\bs{n}_1=(1,0,0)\). 在平面 \(BSC\) 内取 \(\overrightarrow{BC}=(0,-1,0)\)，\(\overrightarrow{BS}=(-1,-1,1)\)，则可取平面 \(BSC\) 的法向量 \(\bs{n}_2=\overrightarrow{BC}\times\overrightarrow{BS}=(-1,0,-1)\). 因此两平面所成的锐二面角 \(\varphi\) 满足 \(\cos\varphi=\dfrac{|\bs{n}_1\cdot\bs{n}_2|}{|\bs{n}_1||\bs{n}_2|}=\dfrac1{\sqrt2}\)，所以 \(\varphi=45^{\circ}\).
\item 因为 \(M\) 是 \(SA\) 的中点，所以 \(M\left(0,\dfrac12,\dfrac12\right)\). 于是 \(\overrightarrow{DM}=\left(0,\dfrac12,\dfrac12\right)\)，\(\overrightarrow{SB}=(1,1,-1)\)，从而 \(\overrightarrow{DM}\cdot\overrightarrow{SB}=0\). 故异面直线 \(DM\) 与 \(SB\) 所成的角为 \(90^{\circ}\).
\end{enumerate}
\end{solution}

\begin{problem}
已知点 \(A(2,8)\)，\(B(x_1,y_1)\)，\(C(x_2,y_2)\) 在抛物线 \(y^2=2px\) 上，\(\triangle ABC\) 的重心与此抛物线的焦点 \(F\) 重合（如图）.

\begin{enumerate}
\item 写出该抛物线的方程和焦点 \(F\) 的坐标；
\item 求线段 \(BC\) 中点 \(M\) 的坐标；
\item 求 \(BC\) 所在直线的方程.
\end{enumerate}

\bitmapinlinefigure[7.6cm][]{img/2004/beijing_science_spring/q18_fig1.png}
\end{problem}

\begin{solution}
\begin{enumerate}
\item 因为点 \(A(2,8)\) 在抛物线 \(y^2=2px\) 上，所以 \(8^2=2p\cdot2\)，解得 \(p=16\). 对于抛物线 \(y^2=2px\)，焦点为 \(F\left(\dfrac p2,0\right)\)，故抛物线方程为 \(y^2=32x\)，焦点为 \(F(8,0)\).
\item 设 \(\triangle ABC\) 的重心为 \(G\). 由 \(G=F(8,0)\)，得 \(\left(\dfrac{2+x_1+x_2}{3},\dfrac{8+y_1+y_2}{3}\right)=(8,0)\)，所以 \(x_1+x_2=22\)，\(y_1+y_2=-8\). 因此线段 \(BC\) 的中点为 \(M\left(\dfrac{x_1+x_2}{2},\dfrac{y_1+y_2}{2}\right)=(11,-4)\).
\item 由重心条件已得 \(y_1+y_2=-8\). 若 \(x_1=x_2\)，则由 \(B\)，\(C\) 在抛物线 \(y^2=32x\) 上，有 \(y_1^2=y_2^2\)；由于 \(B\ne C\)，只能有 \(y_1=-y_2\)，这与 \(y_1+y_2=-8\) 矛盾. 因此 \(BC\) 不是竖直直线，可设其方程为 \(y=kx+b\). 将 \(x=\dfrac{y-b}{k}\) 代入抛物线方程 \(y^2=32x\)，得 \(y^2-\dfrac{32}{k}y+\dfrac{32b}{k}=0\)，其两根为 \(y_1\)，\(y_2\). 由韦达定理，\(y_1+y_2=\dfrac{32}{k}=-8\)，故 \(k=-4\). 又因为 \(B\)，\(C\) 在直线 \(BC\) 上，故 \(y_1+y_2=k(x_1+x_2)+2b\)，即 \(-8=-4\times22+2b\)，解得 \(b=40\). 所以 \(BC\) 所在直线的方程为 \(y=-4x+40\)，即 \(4x+y-40=0\).
\end{enumerate}
\end{solution}

\begin{problem}
某厂生产某种零件，每个零件的成本为 \(40\text{ 元}\)，出厂单价定为 \(60\text{ 元}\). 该厂为鼓励销售商订购，决定当一次订购量超过 \(100\text{ 个}\) 时，每多订购 \(1\text{ 个}\)，订购的全部零件的出厂单价就降低 \(0.02\text{ 元}\)，但实际出厂单价不能低于 \(51\text{ 元}\).

\begin{enumerate}
\item 当一次订购量为多少\(\text{个}\)时，零件的实际出厂单价恰降为 \(51\text{ 元}\)？
\item 设一次订购量为 \(x\text{ 个}\)，零件的实际出厂单价为 \(P\text{ 元}\)，写出函数 \(P=f(x)\) 的表达式；
\item 当销售商一次订购 \(500\text{ 个}\)零件时，该厂获得的利润是多少\(\text{元}\)？如果订购 \(1000\text{ 个}\)，利润又是多少\(\text{元}\)？（工厂售出一个零件的利润 \(=\) 实际出厂单价 \(-\) 成本）
\end{enumerate}
\end{problem}

\begin{solution}
\begin{enumerate}
\item 当 \(x>100\) 时，按规则计算得到的单价为 \(60-0.02(x-100)=62-0.02x\). 令实际单价降为 \(51\text{ 元}\)，得 \(62-0.02x=51\)，解得 \(x=550\). 因此一次订购 \(550\text{ 个}\)时，实际出厂单价恰为 \(51\text{ 元}\).
\item 当 \(0<x\le100\) 时，单价仍为 \(60\text{ 元}\)；当 \(100<x<550\) 时，单价为 \(62-0.02x\text{ 元}\)；当 \(x\ge550\) 时，按规定单价不再降低，保持为 \(51\text{ 元}\). 所以
\[
P=f(x)=
\begin{cases}
60, & 0<x\le100,\\
62-0.02x, & 100<x<550,\\
51, & x\ge550.
\end{cases}
\]
\item 订购 \(500\text{ 个}\)时，单价为 \(62-0.02\times500=52\text{ 元}\)，每个零件的利润为 \(52-40=12\text{ 元}\)，所以总利润为 \(500\times12=6000\text{ 元}\). 订购 \(1000\text{ 个}\)时，单价为 \(51\text{ 元}\)，每个零件的利润为 \(51-40=11\text{ 元}\)，所以总利润为 \(1000\times11=11000\text{ 元}\).
\end{enumerate}
\end{solution}

\begin{problem}
下表给出一个“等差数阵”：

\begin{center}
\small
\begin{tabular}{|c|c|c|c|c|c|c|c|}
\hline
\(4\) & \(7\) & \(\left(\ \right)\) & \(\left(\ \right)\) & \(\left(\ \right)\) & \(\cdots\) & \(a_{1j}\) & \(\cdots\) \\
\hline
\(7\) & \(12\) & \(\left(\ \right)\) & \(\left(\ \right)\) & \(\left(\ \right)\) & \(\cdots\) & \(a_{2j}\) & \(\cdots\) \\
\hline
\(\left(\ \right)\) & \(\left(\ \right)\) & \(\left(\ \right)\) & \(\left(\ \right)\) & \(\left(\ \right)\) & \(\cdots\) & \(a_{3j}\) & \(\cdots\) \\
\hline
\(\left(\ \right)\) & \(\left(\ \right)\) & \(\left(\ \right)\) & \(\left(\ \right)\) & \(\left(\ \right)\) & \(\cdots\) & \(a_{4j}\) & \(\cdots\) \\
\hline
\(\vdots\) & \(\vdots\) & \(\vdots\) & \(\vdots\) & \(\vdots\) & \(\cdots\) & \(\vdots\) & \(\cdots\) \\
\hline
\(a_{i1}\) & \(a_{i2}\) & \(a_{i3}\) & \(a_{i4}\) & \(a_{i5}\) & \(\cdots\) & \(a_{ij}\) & \(\cdots\) \\
\hline
\(\vdots\) & \(\vdots\) & \(\vdots\) & \(\vdots\) & \(\vdots\) & \(\cdots\) & \(\vdots\) & \(\cdots\) \\
\hline
\end{tabular}
\end{center}

其中每行，每列都是等差数列，\(a_{ij}\) 表示位于第 \(i\) 行第 \(j\) 列的数.

\begin{enumerate}
\item 写出 \(a_{45}\) 的值；
\item 写出 \(a_{ij}\) 的计算公式；
\item 证明：正整数 \(N\) 在该等差数列阵中的充要条件是 \(2N+1\) 可以分解成两个不是 \(1\) 的正整数之积.
\end{enumerate}
\end{problem}

\begin{solution}
\begin{enumerate}
\item 第 \(i\) 行的公差记为 \(r_i\)，第 \(j\) 列的公差记为 \(s_j\). 由第一行的前两项和第二行的前两项，得 \(r_1=7-4=3\)，\(r_2=12-7=5\). 对任意 \(i,j\)，比较相邻两个 \(2\times2\) 小方格的两条路径，得 \(r_{i+1}-r_i=s_{j+1}-s_j\). 左端只与 \(i\) 有关，右端只与 \(j\) 有关，因此二者都等于同一个常数. 取 \(i=1\)，由 \(r_2-r_1=2\) 得 \(s_{j+1}-s_j=2\)；再取 \(j=1\)，得 \(r_{i+1}-r_i=2\). 因而 \(r_i=3+2(i-1)=2i+1\). 同理，由 \(s_1=7-4=3\) 得 \(s_j=3+2(j-1)=2j+1\).
\item 第一列的公差为 \(s_1=3\)，所以 \(a_{i1}=4+3(i-1)=3i+1\). 沿第 \(i\) 行递推，得 \(a_{ij}=a_{i1}+(j-1)r_i=3i+1+(j-1)(2i+1)=2ij+i+j\). 因此 \(a_{45}=2\times4\times5+4+5=49\).
\item 对任意正整数 \(i,j\)，由上式有 \(a_{ij}=2ij+i+j\)，所以 \(2a_{ij}+1=4ij+2i+2j+1=(2i+1)(2j+1)\). 两个因数 \(2i+1\)，\(2j+1\) 都是大于 \(1\) 的正整数，故若正整数 \(N\) 在数阵中，则 \(2N+1\) 可以分解成两个不是 \(1\) 的正整数之积.

反之，若正整数 \(N\) 满足 \(2N+1=uv\)，其中 \(u,v\) 是大于 \(1\) 的正整数，则 \(2N+1\) 为奇数，故 \(u,v\) 都是奇数. 于是存在正整数 \(i,j\)，使得 \(u=2i+1\)，\(v=2j+1\). 因此 \(2N+1=(2i+1)(2j+1)=4ij+2i+2j+1\)，从而 \(N=2ij+i+j=a_{ij}\)，即 \(N\) 在该等差数阵中. 综上，所述条件是充要的.
\end{enumerate}
\end{solution}
